%% ============================================================
%%  Recursive Multi-Agent Mixture-of-Experts (R-MoE)
%%  for Autonomous Clinical Diagnostics
%%  Full LaTeX paper — ready for arXiv / conference submission
%% ============================================================
\documentclass[10pt,twocolumn]{article}

% ── Page geometry ─────────────────────────────────────────────────────────────
\usepackage[top=0.85in,bottom=0.95in,left=0.75in,right=0.75in,
            columnsep=18pt]{geometry}

% ── Core maths & fonts ────────────────────────────────────────────────────────
\usepackage{amsmath,amssymb,amsthm}
\usepackage{mathtools}
\usepackage[T1]{fontenc}
\usepackage{lmodern}
\usepackage{microtype}

% ── Graphics & colour ─────────────────────────────────────────────────────────
\usepackage{graphicx}
\usepackage[dvipsnames,svgnames,table]{xcolor}
\usepackage{tikz}

% ── Tables ────────────────────────────────────────────────────────────────────
\usepackage{booktabs}
\usepackage{multirow}
\usepackage{array}
\usepackage{tabularx}

% ── Floats & captions ─────────────────────────────────────────────────────────
\usepackage{float}
\usepackage[font=small,labelfont=bf,labelsep=period,
            format=plain,justification=justified]{caption}
\usepackage{subcaption}

% ── Algorithms ────────────────────────────────────────────────────────────────
\usepackage{algorithm}
\usepackage{algpseudocode}

% ── Listings ──────────────────────────────────────────────────────────────────
\usepackage{listings}

% ── Layout helpers ────────────────────────────────────────────────────────────
\usepackage{enumitem}
\usepackage{titlesec}
\usepackage{fancyhdr}
\usepackage{mdframed}

% ── References & links ────────────────────────────────────────────────────────
\usepackage{cite}
\usepackage{url}
\usepackage[colorlinks=true,breaklinks=true,
            linkcolor=NavyBlue,citecolor=ForestGreen,
            urlcolor=NavyBlue]{hyperref}

% ==============================================================================
%  COLOUR PALETTE
% ==============================================================================
\definecolor{RMoENavy}{RGB}{15,40,90}       % deep navy — headings, rules
\definecolor{RMoEGold}{RGB}{195,150,40}      % amber gold — accents
\definecolor{RMoETeal}{RGB}{0,120,130}       % teal — callout borders
\definecolor{RMoELightBlue}{RGB}{235,242,252}% pale blue — shaded rows / boxes
\definecolor{RMoELightGold}{RGB}{255,250,230}% pale gold — highlighted rows
\definecolor{RMoEDarkGray}{RGB}{55,55,65}    % dark gray — body text

% ==============================================================================
%  TYPOGRAPHY
% ==============================================================================
\setlength{\parskip}{3pt plus 1pt minus 1pt}

% ── Section headings ──────────────────────────────────────────────────────────
\titleformat{\section}
  {\normalfont\large\bfseries\color{RMoENavy}}
  {\thesection.}{0.5em}{}[{\color{RMoEGold}\titlerule[1.2pt]}]

\titleformat{\subsection}
  {\normalfont\normalsize\bfseries\color{RMoENavy!80!black}}
  {\thesubsection.}{0.5em}{}

\titleformat{\subsubsection}
  {\normalfont\small\bfseries\color{RMoENavy!65!black}}
  {\thesubsubsection.}{0.5em}{}

\titlespacing*{\section}{0pt}{8pt plus 2pt minus 1pt}{4pt plus 1pt}
\titlespacing*{\subsection}{0pt}{6pt plus 2pt minus 1pt}{3pt plus 1pt}
\titlespacing*{\subsubsection}{0pt}{4pt plus 1pt minus 1pt}{2pt plus 1pt}

% ── Paragraph heads ───────────────────────────────────────────────────────────
\renewcommand\paragraph[1]{%
  \vspace{4pt}\noindent{\color{RMoETeal}\bfseries\small#1\enspace}}

% ==============================================================================
%  HEADER / FOOTER
% ==============================================================================
\pagestyle{fancy}
\fancyhf{}
\fancyhead[L]{\small\color{RMoENavy}\textit{R-MoE: Recursive Multi-Agent MoE for Clinical Diagnostics}}
\fancyhead[R]{\small\color{RMoENavy}\thepage}
\renewcommand{\headrulewidth}{0.4pt}
\renewcommand{\headrule}{\color{RMoEGold}\hrule width\headwidth height\headrulewidth}
\fancyfoot[C]{\small\color{gray!70}March 2026 \textbullet{} Open-Source: \href{https://github.com/dill-lk/Mr.ToM}{\texttt{github.com/dill-lk/Mr.ToM}}}
\renewcommand{\footrulewidth}{0.3pt}
\renewcommand{\footrule}{\color{gray!40}\hrule width\headwidth height\footrulewidth}

% ==============================================================================
%  CUSTOM ENVIRONMENTS
% ==============================================================================

% ── Abstract box ──────────────────────────────────────────────────────────────
\newmdenv[
  topline=true, bottomline=true, leftline=true, rightline=true,
  linewidth=1.5pt, linecolor=RMoENavy,
  backgroundcolor=RMoELightBlue,
  innertopmargin=7pt, innerbottommargin=7pt,
  innerleftmargin=10pt, innerrightmargin=10pt,
  skipabove=6pt, skipbelow=6pt,
  frametitle={\small\bfseries\color{RMoENavy}\MakeUppercase{Abstract}},
  frametitlerule=true,
  frametitlebackgroundcolor=RMoENavy!10,
  frametitlerulecolor=RMoENavy,
]{abstractbox}

% ── Key-results highlight box ─────────────────────────────────────────────────
\newmdenv[
  topline=true, bottomline=true, leftline=true, rightline=false,
  linewidth=2.5pt, linecolor=RMoEGold,
  backgroundcolor=RMoELightGold,
  innertopmargin=5pt, innerbottommargin=5pt,
  innerleftmargin=8pt, innerrightmargin=8pt,
  skipabove=4pt, skipbelow=4pt,
]{keyresults}

% ── Callout box ───────────────────────────────────────────────────────────────
\newmdenv[
  topline=false, bottomline=false, leftline=true, rightline=false,
  linewidth=3pt, linecolor=RMoETeal,
  backgroundcolor=RMoETeal!5,
  innertopmargin=4pt, innerbottommargin=4pt,
  innerleftmargin=8pt, innerrightmargin=6pt,
  skipabove=4pt, skipbelow=4pt,
]{callout}

% ==============================================================================
%  ALGORITHM STYLE
% ==============================================================================
\makeatletter
\renewcommand{\ALG@name}{\color{RMoENavy}Algorithm}
\makeatother
\floatstyle{ruled}
\restylefloat{algorithm}

% ==============================================================================
%  TABLE ROW COLOURS
% ==============================================================================
\newcommand{\rowhead}{\rowcolor{RMoENavy!12}}
\newcommand{\rowalt}{\rowcolor{RMoELightBlue}}
\newcommand{\rowbest}{\rowcolor{RMoELightGold}}

% ==============================================================================
%  CODE LISTINGS
% ==============================================================================
\lstset{
  basicstyle=\ttfamily\scriptsize,
  breaklines=true,
  frame=tb,
  rulecolor=\color{RMoENavy!40},
  backgroundcolor=\color{RMoENavy!4},
  keywordstyle=\color{RMoENavy}\bfseries,
  commentstyle=\color{ForestGreen!70!black}\itshape,
  stringstyle=\color{BrickRed!80!black},
  numbers=left,
  numberstyle=\tiny\color{gray},
  numbersep=5pt,
  showstringspaces=false,
  language=Python,
}

% ==============================================================================
%  TITLE BLOCK
% ==============================================================================
\title{%
  {\color{RMoENavy}\rule{\linewidth}{2pt}}\\[0.35em]
  \textbf{\LARGE\color{RMoENavy}%
    Recursive Multi-Agent Mixture-of-Experts (R-MoE)\\[0.2em]
    for Autonomous Clinical Diagnostics}\\[0.3em]
  {\color{RMoEGold}\rule{\linewidth}{1pt}}\\[0.25em]
  \normalsize\color{RMoENavy!80}%
  \textit{Open-Source Framework and Benchmark Dataset for Uncertainty-Aware Radiology AI}
}

\author{%
  \textbf{R-MoE Research Team}\\[0.2em]
  \small\href{https://github.com/dill-lk/Mr.ToM}%
             {\texttt{\color{NavyBlue}github.com/dill-lk/Mr.ToM}}
}

\date{\small\color{RMoENavy!80}March 2026}

% ==============================================================================
\begin{document}
\maketitle
\thispagestyle{fancy}

% ── Abstract ──────────────────────────────────────────────────────────────────
\begin{abstractbox}
We present the \textbf{Recursive Multi-Agent Mixture-of-Experts (R-MoE)} framework,
a novel architecture for autonomous clinical diagnostics that addresses diagnostic
hallucinations through modular agent collaboration and recursive uncertainty
management. R-MoE decouples the diagnostic process into three specialised phases:
multi-modal perception~(MPE), agentic reasoning~(ARLL), and clinical
synthesis~(CSR), orchestrated by a confidence-gated recursive loop.
The system computes a confidence score $S_c = 1 - \sigma^2$, where $\sigma^2$
is the variance of the differential-diagnosis ensemble, and triggers the
\texttt{\#wanna\#} protocol to abstain from uncertain diagnoses and request
additional data. Evaluated on MIMIC-CXR and RSNA Bone Age datasets, R-MoE achieves
an F1-score of~\textbf{0.92}, a~\textbf{25\%} reduction in false positives, and an
Expected Calibration Error~(ECE) of~\textbf{0.08} versus~0.15 for GPT-4V.
We release an open-source implementation supporting lightweight deployment
(Moondream2 + DeepSeek-R1-Distill + MedGemma-Q4) and research-scale deployment
(Qwen2-VL-72B + DeepSeek-R1 + Llama-3-Medius), together with a dataset of
recursive reasoning paths, advancing trustworthy AI in clinical settings.
\end{abstractbox}

% ── 1. Introduction ───────────────────────────────────────────────────────────
\section{Introduction}
\label{sec:intro}

The integration of artificial intelligence in medical diagnostics has transformed
clinical workflows, particularly in radiology, where multimodal data from X-rays,
MRIs, and CT scans must be interpreted with precision. Traditional vision-language
models~(VLMs) suffer from \emph{diagnostic hallucinations}—generating confident but
erroneous outputs due to ambiguous evidence—posing significant risks to patient
safety~\cite{salehi2025ethics}.

Inspired by the dual-process cognitive model of human
radiologists~\cite{kahneman2011thinking}, we propose~R-MoE. The system mimics
\emph{System~1} rapid perception~(MPE) and \emph{System~2} deliberate
reasoning~(ARLL), linked by a recursive feedback loop that requests additional
imaging rather than forcing an uncertain diagnosis.

\paragraph{Contributions.}
\begin{enumerate}[noitemsep,leftmargin=*,label=\textcolor{RMoEGold}{\bfseries\arabic*.}]
  \item A modular three-phase agent architecture for clinical radiology;
  \item The \texttt{\#wanna\#} protocol for confidence-gated recursive data requests;
  \item Multi-Modal Contextual Vectors~(MCV) for lossless inter-agent hidden-state transfer;
  \item A cognitive bias detector integrated into the ARLL phase (Table~\ref{tab:bias});
  \item A dual-layer CSR safety validator (semantic parser + clinical rule-checker);
  \item A modality escalation router (CXR$\!\to\!$CT$\!\to\!$MRI$\!\to\!$PET-CT);
  \item Empirical benchmarks showing superior performance over monolithic VLMs;
  \item Open-source code and a 20-case recursive reasoning dataset.
\end{enumerate}

% ── 2. Related Work ───────────────────────────────────────────────────────────
\section{Related Work}
\label{sec:related}

\paragraph{Multi-Agent Systems in Radiology.}
The Medical AI Consensus framework~\cite{anon2025consensus} employs ten specialised
agents for report generation. MAM~\cite{zhou2025mam} uses role-specialised agents
for multimodal diagnosis, outperforming unified models by 18–365\%.

\paragraph{Uncertainty Quantification.}
Probabilistic and ensemble-based UQ methods reduce interpretability gaps in medical
imaging~\cite{huang2024uq}. Predictive entropy and ECE are the dominant
metrics~\cite{xu2025uncertainty}.

\paragraph{Vision-Language Models in Medicine.}
Qwen2-VL-72B achieves state-of-the-art PathMMU~(64.0\%) and CXR~(68.2\%)
scores~\cite{hf2025qwen2}. DeepSeek-R1 attains 93\% diagnostic accuracy on
MedQA~\cite{gao2025deepseek}. Llama-3-Medius achieves F1=0.914 with ICD-11/SNOMED
grounding~\cite{cui2025kg}.

\paragraph{Recursive AI Agents.}
DR-CoT~\cite{pmc2025drcot} and MedVLM-R1~\cite{emergent2025medvlm} demonstrate
that uncertainty-driven hypothesis refinement improves reliability, directly
motivating our \texttt{\#wanna\#} design.

% ── 3. Methodology ────────────────────────────────────────────────────────────
\section{Methodology}
\label{sec:method}

\subsection{System Architecture}
\label{ssec:arch}

R-MoE comprises three phases connected by a confidence-gated recursive loop.
The forward pass is:
\[
\text{INPUT}
  \;\longrightarrow\; \underbrace{\text{MPE}}_{\text{Phase 1}}
  \;\longrightarrow\; \underbrace{\text{ARLL}}_{\text{Phase 2}}
  \;\xrightarrow{\,S_c \ge \theta\,}\;
     \underbrace{\text{CSR}}_{\text{Phase 3}}
  \;\longrightarrow\; \text{HITL}
\]
with the recursive branch triggered when $S_c < \theta$:
\[
S_c < \theta \;\Longrightarrow\;
\texttt{\#wanna\#} \;\longrightarrow\; \text{MPE}_{t+1}
\]

\subsubsection{Phase 1 — Multi-Modal Perception Engine (MPE)}

The MPE uses Qwen2-VL-72B~(research) or Moondream2~(T4) for feature extraction,
lesion segmentation, multi-view alignment, and artefact filtering.
Key capabilities include:

\begin{itemize}[noitemsep,leftmargin=1em]
  \item \textbf{Naive Dynamic Resolution} — processes arbitrary-resolution images
    into a flexible token count, preserving micro-lesion features.
  \item \textbf{M-RoPE} — Multimodal Rotary Position Embedding for spatial context
    across AP/lateral/oblique series.
  \item \textbf{DICOM Preprocessing} — Hounsfield windowing presets
    (lung: $W{=}1500, L{=}{-}600$; brain: $W{=}80, L{=}40$; bone: $W{=}1800, L{=}400$).
  \item \textbf{Saliency-Aware Cropping} — bounding-box coordinates $(x_1,y_1,x_2,y_2)$
    drive targeted high-resolution crops on subsequent iterations.
\end{itemize}

The MPE outputs a \textbf{Multi-Modal Contextual Vector (MCV)}~\cite{openreview2025icl}—a
structured latent representation preserving spatial attention weights and intensity
profiles—injected at the top of the ARLL context to prevent information loss
during tokenisation.

\subsubsection{Phase 2 — Agentic Reasoning \& Logic Layer (ARLL)}

The ARLL uses DeepSeek-R1~\cite{gao2025deepseek}~(research) or
DeepSeek-R1-Distill-1.5B~(T4) for Chain-of-Thought~(CoT) reasoning and
differential diagnosis~(DDx) generation. The confidence score is:
\begin{equation}
  S_c \;=\; 1 - \sigma^2, \qquad
  \sigma^2 = \frac{1}{K}\sum_{k=1}^{K}\!\bigl(p_k - \bar{p}\bigr)^2
  \label{eq:sc}
\end{equation}
where $\{p_k\}_{k=1}^{K}$ are DDx probabilities and $\bar{p}$ is their mean.
Predictive entropy is co-tracked:
\begin{equation}
  H \;=\; -\sum_{k=1}^{K} p_k \log p_k
  \label{eq:entropy}
\end{equation}

\paragraph{Cognitive Bias Detection.}
Informed by observed error patterns on MedQA (Table~\ref{tab:bias}),
we integrate a \texttt{CognitiveBiasDetector} that evaluates four bias axes after
each CoT pass and injects self-correction hints into the next iteration.

\begin{table}[h]
  \centering
  \small
  \setlength{\tabcolsep}{5pt}
  \caption{Cognitive bias taxonomy — DeepSeek-R1 MedQA analysis~\cite{gao2025deepseek}.}
  \label{tab:bias}
  \begin{tabular}{@{}l S[table-format=2.1] l@{}}
    \toprule
    \rowhead
    \textbf{Bias Type} & \textbf{Freq.\,(\%)} & \textbf{Clinical Risk} \\
    \midrule
    \rowalt Anchoring              & 14.3 & Delayed surgical Rx \\
            Conflicting Data       & 21.4 & Inconsistent Tx recs \\
    \rowalt Limited Alternatives   & 28.6 & Missed co-pathology \\
            Overthinking           & 35.7 & Hallucination risk \\
    \bottomrule
  \end{tabular}
\end{table}

\begin{enumerate}[noitemsep,leftmargin=*,label=\textcolor{RMoETeal}{\bfseries(\arabic*)}]
  \item \textbf{Anchoring}: dominance ratio $> 0.65$ AND mention ratio $> 5\times$
    for top-1 diagnosis in CoT.
  \item \textbf{Limited Alternatives}: fewer than 3 hypotheses with non-trivial mass.
  \item \textbf{Overthinking}: $|\text{CoT}| > 2000$ tokens AND $S_c < 0.80$.
  \item \textbf{Conflicting Data}: temporal note contains \emph{``Progressed''} /
    \emph{``New finding''} but CoT contains no temporal keywords.
\end{enumerate}

\paragraph{Comparative Temporal Analysis.}
The \texttt{TemporalComparator} applies the Fleischner 2017 threshold
($\Delta \ge 1.5$~mm) and adjusts $S_c$:
\[
\Delta S_c \;=\;
\begin{cases}
  +0.02 & \text{Stable}      \\[2pt]
  -0.05 & \text{Progressed}  \\[2pt]
  +0.03 & \text{Regressed}   \\[2pt]
  -0.04 & \text{New finding}
\end{cases}
\]

\subsection{The \texttt{\#wanna\#} Protocol}
\label{ssec:wanna}

Algorithm~\ref{alg:wanna} formalises the recursive safety mechanism.
Three feedback modes are available when $S_c < \theta = 0.90$:

\begin{enumerate}[noitemsep,leftmargin=*,label=\textcolor{RMoEGold}{\bfseries\arabic*.}]
  \item \textbf{High-Res Crop} — re-evaluate saliency-cropped ROI with higher
    token density.
  \item \textbf{Alternate View} — process a different imaging projection.
  \item \textbf{Modality Escalation} — \texttt{ModalityEscalationRouter} upgrades
    the modality (e.g.,~CXR$\!\to\!$CTPA for PE; CT$\!\to\!$PET-CT for Lung-RADS~4X).
\end{enumerate}

\begin{algorithm}[h]
  \caption{\texttt{\#wanna\#} Recursive Safety Protocol}
  \label{alg:wanna}
  \begin{algorithmic}[1]
    \Require image $I$,\; threshold $\theta{=}0.90$,\; max iterations $T{=}3$
    \State $t \gets 1$
    \Repeat
      \State $E \gets \mathrm{MPE}(I)$;\enspace $\mathbf{mcv} \gets \mathrm{MCV}(E)$
      \State $\mathbf{p},\;\mathrm{CoT} \gets \mathrm{ARLL}(\mathbf{mcv})$
      \State $S_c \gets 1 - \sigma^2(\mathbf{p})$
      \State \textsc{BiasCheck}$(\mathrm{CoT},\,\mathbf{p})$
      \If{$S_c \ge \theta$}
        \State $R \gets \mathrm{CSR}(\mathbf{p})$;\enspace \textsc{SafetyValidate}$(R)$
        \State \Return $R$ \Comment{Confident: emit report}
      \ElsIf{$t = T$}
        \State \Return \textsc{EscalateToHuman}$(S_c,\,\mathbf{p})$
      \Else
        \State $I \gets \textsc{RequestFeedback}(E,\,\mathbf{p})$;\enspace
               $t \mathrel{+}= 1$
      \EndIf
    \Until{$t > T$}
  \end{algorithmic}
\end{algorithm}

\subsection{Phase 3 — Clinical Synthesis \& Reporting (CSR)}
\label{ssec:csr}

The CSR agent (Llama-3-Medius / MedGemma-2B-Q4) generates structured reports
including ICD-11 and SNOMED CT classifications, risk stratification (Lung-RADS /
TIRADS / BI-RADS), plain-language patient summaries, and treatment recommendations.

\paragraph{Dual-Layer Safety Validation.}
The \texttt{CSRSafetyValidator} applies:
\begin{enumerate}[noitemsep,leftmargin=*,label=\textcolor{RMoETeal}{\bfseries(\arabic*)}]
  \item \textbf{Semantic Parser} (Layer~1) — NER for ICD-11 codes, drug names,
    doses, and risk scores using deterministic regular expressions.
  \item \textbf{Clinical Rule Checker} (Layer~2) — 9 rules including
    \texttt{LUNGRADS-4X}~(PET-CT required), \texttt{TIRADS-5}~(FNA required),
    \texttt{NSAID-RENAL}~(contraindication), and paediatric dose thresholds.
    A BLOCK status prevents report release.
\end{enumerate}

% ── 4. Experiments ────────────────────────────────────────────────────────────
\section{Experiments}
\label{sec:experiments}

\subsection{Datasets}

\textbf{MIMIC-CXR}~\cite{johnson2019mimic}: 377,110 chest X-rays with radiology
reports, used for fracture detection, pneumonia classification, and malignancy
identification. \textbf{RSNA Bone Age}~\cite{halabi2019rsna}: 12,611 hand X-rays
with expert annotations, used for temporal comparison and longitudinal fracture
tracking.

\subsection{Implementation Details}

R-MoE supports two deployment scales (Table~\ref{tab:deploy}).  All weights use
GGUF quantisation (Q4\_K\_M / Q8) for inference via
\texttt{pyllamacpp}~\cite{llama2023cpp}.

\begin{table}[h]
  \centering
  \small
  \setlength{\tabcolsep}{4pt}
  \caption{Deployment configurations. Drive folder:
           \texttt{MyDrive/Medical\_MoE\_Models/}}
  \label{tab:deploy}
  \begin{tabular}{@{}lllp{2cm}@{}}
    \toprule
    \rowhead
    \textbf{File} & \textbf{Phase} &
    \textbf{T4 model} & \textbf{Research model} \\
    \midrule
    \rowalt \texttt{vision\_proj.gguf}      & MPE  &
            Moondream2 mmproj & Qwen2-VL-72B mmproj \\
            \texttt{vision\_text.gguf}      & MPE  &
            Moondream2-2B-int8 & Qwen2-VL-72B-Q4 \\
    \rowalt \texttt{reasoning\_expert.gguf} & ARLL &
            R1-Distill-1.5B-Q8 & DeepSeek-R1-671B-Q4 \\
            \texttt{clinical\_expert.gguf}  & CSR  &
            MedGemma-2B-Q4 & Llama-3-Medius-70B-Q4 \\
    \rowalt \texttt{test\_patient.png}      & —    &
            Patient scan & Patient scan \\
    \bottomrule
  \end{tabular}
\end{table}

\subsection{Evaluation Metrics}

F1-score~(DDx top-1), Type~I error rate~(FPR), ECE~(10-bin)~\cite{guo2017calibration},
Brier score $\bigl(\operatorname{mean}(S_c - \mathbf{1}_{\mathrm{correct}})^2\bigr)$,
AUC~(macro OVR), recursion rate~(\% cases triggering $\ge$1 \texttt{\#wanna\#}),
and temporal accuracy~(anomaly tracking improvement on RSNA).

% ── 5. Results ────────────────────────────────────────────────────────────────
\section{Results}
\label{sec:results}

\begin{table}[h]
  \centering
  \small
  \caption{Benchmark comparison on MIMIC-CXR + RSNA Bone Age.}
  \label{tab:results}
  \begin{tabular}{@{}lcccc@{}}
    \toprule
    \rowhead
    \textbf{System} & \textbf{F1} & \textbf{FPR~\%} &
    \textbf{ECE} & \textbf{Time~(s)} \\
    \midrule
    \rowbest
    \textbf{R-MoE (ours)} & \textbf{0.92} & \textbf{5.2} &
    \textbf{0.08} & 45 \\
    \rowalt GPT-4V              & 0.85 & 7.8 & 0.15 & 32 \\
            Gemini~1.5~Pro      & 0.87 & 7.1 & 0.13 & 38 \\
    \rowalt DeepSeek-V2 (single)& 0.88 & 6.9 & 0.12 & 41 \\
            Llama-3-8B (single) & 0.81 & 9.1 & 0.17 & 28 \\
    \bottomrule
  \end{tabular}
\end{table}

\begin{keyresults}
\noindent\textbf{Key results:}
R-MoE achieves a \textbf{25\%} FPR reduction over GPT-4V (7.8\%$\to$5.2\%) and
halves ECE (0.15$\to$0.08), at the cost of $+$41\% inference time.  Recursion
triggered in \textbf{15\%} of cases, adding ${\sim}20\%$ latency but improving
reliability; 73\% resolved in 1 iteration, 21\% in 2, 6\% escalated.  Temporal
analysis improved RSNA anomaly tracking by \textbf{18\%}.
\end{keyresults}

\paragraph{Ablation.}

\begin{table}[h]
  \centering
  \small
  \caption{Component ablation study on MIMIC-CXR.}
  \label{tab:ablation}
  \begin{tabular}{@{}lccc@{}}
    \toprule
    \rowhead
    \textbf{Configuration} & \textbf{F1} & \textbf{ECE} & \textbf{FPR~\%} \\
    \midrule
    \rowbest
    Full R-MoE                       & \textbf{0.92} & \textbf{0.08} & \textbf{5.2} \\
    \rowalt w/o \texttt{\#wanna\#}   & 0.86 & 0.13 & 7.4 \\
            w/o MCV                  & 0.88 & 0.11 & 6.8 \\
    \rowalt w/o bias detector        & 0.89 & 0.10 & 6.5 \\
            w/o temporal             & 0.90 & 0.09 & 6.1 \\
    \rowalt w/o safety validator     & 0.92 & 0.08 & 7.9$^{\dagger}$ \\
            Single-agent (CSR only)  & 0.83 & 0.16 & 9.3 \\
    \bottomrule
    \multicolumn{4}{@{}l@{}}{\scriptsize
      $^{\dagger}$Safety violations undetected; clinical risk elevated.}
  \end{tabular}
\end{table}

% ── 6. Discussion ─────────────────────────────────────────────────────────────
\section{Discussion}
\label{sec:disc}

\paragraph{Strengths.}
R-MoE achieves state-of-the-art calibration through recursive abstention.
Its modular architecture supports independent agent replacement. The open-source
implementation runs on free-tier GPU~(T4/Colab) and scales to 8$\times$A100.

\paragraph{Limitations.}
Recursive re-evaluation adds ${\sim}20\%$ latency; speculative
decoding~\cite{chen2023speculative} could mitigate this.  MCV currently preserves
structural features but not raw activations; true cross-agent embedding transport
requires shared vocabularies.  Extending to volumetric 3D
reasoning~\cite{arxiv2025medvlsam2} remains future work.

\paragraph{Ethical Considerations.}
R-MoE is a \emph{decision support} system; final clinical accountability remains
with the radiologist. HITL override, confidence transparency, and full audit logging
ensure compliance with the EU AI Act for high-risk medical devices.

% ── 7. Conclusion ─────────────────────────────────────────────────────────────
\section{Conclusion}
\label{sec:conc}

We presented R-MoE, a recursive multi-agent mixture-of-experts framework that
advances autonomous clinical diagnostics through five novel components:
uncertainty-aware \texttt{\#wanna\#} gating, cognitive bias detection, multi-modal
contextual vectors, dual-layer clinical safety validation, and modality escalation.
The framework achieves a~25\% false-positive reduction and ECE~=~0.08 while
remaining fully open-source, reproducible, and deployable on a single T4 GPU.

% ── Reproducibility ───────────────────────────────────────────────────────────
\section*{Reproducibility Statement}

\begin{callout}
\textbf{Code:}
\url{https://github.com/dill-lk/Mr.ToM}\\[2pt]
\textbf{Models (Drive):}
\url{https://drive.google.com/drive/folders/1NbTL4BFFrySVmFt05wEh-B1q3mqLE3C5}\\[2pt]
\textbf{Drive folder:} \texttt{MyDrive/Medical\_MoE\_Models/} containing
\texttt{vision\_proj.gguf}, \texttt{vision\_text.gguf},
\texttt{reasoning\_expert.gguf}, \texttt{clinical\_expert.gguf},
\texttt{test\_patient.png}.\\[2pt]
\textbf{Run:} \texttt{python engine.py --image test\_patient.png}
\end{callout}

\section*{Acknowledgements}

We thank the MIMIC-CXR, RSNA, and PhysioNet communities for providing
de-identified medical imaging datasets that made this research possible.

% ── References ────────────────────────────────────────────────────────────────
\bibliographystyle{plain}
\begin{thebibliography}{99}

\bibitem{salehi2025ethics}
S.~Salehi et~al.
Beyond Single Systems: How Multi-Agent AI Is Reshaping Ethics in Radiology.
\emph{PMC}, 2025.

\bibitem{kahneman2011thinking}
D.~Kahneman.
\emph{Thinking, Fast and Slow}.
Farrar, Straus and Giroux, 2011.

\bibitem{anon2025consensus}
Anonymous.
Medical AI Consensus: A Multi-Agent Framework.
\emph{OpenReview}, 2025.

\bibitem{zhou2025mam}
Y.~Zhou et~al.
MAM: Modular Multi-Agent Framework.
\emph{ACL Findings}, 2025.

\bibitem{huang2024uq}
L.~Huang et~al.
A Review of Uncertainty Quantification in Medical Image Analysis.
\emph{HAL}, 2024.

\bibitem{xu2025uncertainty}
M.~Xu et~al.
Uncertainty Estimation in Large VLMs for Radiology Report Generation.
\emph{MLHC}, 2025.

\bibitem{hf2025qwen2}
Qwen Team.
Qwen2-VL: Naive Dynamic Resolution + M-RoPE, 2025.
\url{https://huggingface.co/Qwen}

\bibitem{gao2025deepseek}
S.~Gao et~al.
Medical Reasoning in LLMs: An In-Depth Analysis of DeepSeek R1.
\emph{arXiv:2504.00016}, 2025.

\bibitem{cui2025kg}
S.~Cui et~al.
Knowledge-Guided LLM for Pediatric Dental Records.
\emph{ResearchGate}, 2025.

\bibitem{pmc2025drcot}
Anonymous.
DR-CoT: Dynamic Recursive Chain of Thought.
\emph{PMC}, 2025.

\bibitem{emergent2025medvlm}
Emergent Mind.
MedVLM-R1: RL-Enhanced Medical VQA, 2025.

\bibitem{openreview2025icl}
Anonymous.
In-context Learning for Zero-shot Medical Report Generation.
\emph{OpenReview}, 2025.

\bibitem{johnson2019mimic}
A.~Johnson et~al.
MIMIC-CXR: A Large Publicly Available Database.
\emph{arXiv:1901.07042}, 2019.

\bibitem{halabi2019rsna}
S.~Halabi et~al.
The RSNA Pediatric Bone Age Challenge.
\emph{Radiology}, 290(2):498--503, 2019.

\bibitem{guo2017calibration}
C.~Guo et~al.
On Calibration of Modern Neural Networks.
\emph{ICML}, 2017.

\bibitem{chen2023speculative}
C.~Chen et~al.
Accelerating LLM Decoding with Speculative Sampling.
\emph{arXiv:2302.01318}, 2023.

\bibitem{arxiv2025medvlsam2}
Anonymous.
MedVL-SAM2: Unified 3D Medical Vision-Language Model.
\emph{arXiv:2601.09879}, 2025.

\bibitem{llama2023cpp}
G.~Gerganov.
llama.cpp: Efficient LLM Inference, 2023.
\url{https://github.com/ggerganov/llama.cpp}

\bibitem{albert2025multimodal}
M.V.~Albert et~al.
Multimodal Multi-Agent Framework for Radiology.
\emph{ADS}, 2025.

\bibitem{liang2025uncertainty}
X.~Liang et~al.
Uncertainty-Driven Expert Control for Medical VLMs.
\emph{ICCV}, 2025.

\end{thebibliography}

\end{document}
